% s-344-p344-m003; original PDF 344, 345
\Aphorism{s-344-p344-m003}{397}
Im Traume revoltirt dasjenige Ich, das vom wachenden überbaut, erstickt, niedergedrückt wird, ein Ich aus früheren tieferen Schichten, das am Tage nicht hervordarf. Der Traum ist Sclavenaufstand und Saturnalienfest; er übt am Wachsein dieselbe subversive Vergeltung wie das jüngste Gericht an der Weltgeschichte: die Ersten werden die Letzten sein. Was am Tage unser Empfinden beherrscht, woran wir hundertmal in einzelnen Bildern oder ununterbrochen in einer Art seelischen „Gemeingefühls“ sinnen und denken, das schweigt im Traume oder wird eigensinnig verschwiegen: kommt es doch einmal hervor, dann vielleicht in einer parodistischen Verzerrung oder in einer regnerischen Alltagsfärbung, dass es zum Erbarmen ist. Umgekehrt, was am Tage unter den Tisch fällt, was unter kräftigem Drucke unten gehalten wird, als das Proletariat der Gedanken und Empfindungen, — Begegnungen, die wir einen Augenblick später schon wieder vergessen haben, Dinge und Menschen, an die wir nicht den Bruchtheil einer Secunde Denkens verschwenden, Begriffskeime, wie sie gleich den organischen Keimen in der Natur zu Millionen auftreten, um zu Millionen unentwickelt wieder abzusterben —: das alles rächt sich für die tagüber erfahrene Unterdrückung und macht sich im Vordergrunde der Traumbühne zu schaffen. Aus den vergessenen Erlebnissen werden drohende Gefahren oder entsetzliche Erinnerungen, aus den übersehenen Menschen hochwichtige Helfer und Widersacher, aus den embryonischen Gedanken fabelhafte Fratzen und Unmöglichkeiten. Lauter falsche Accente, eine verdrehte Rangordnung! Der Traum vergegenwärtigt uns im atavistischen Bilde die Zeiten der beginnenden Cultur, als das Bewusstsein noch ohne Sieb arbeitete, alles einsog und nichts ausschied, als jede Vibration der Gehirnfasern ein neues Weltbild bedeutete, das sich noch keinem erkennbaren Mittelwerthe asymptotisch näherte, — jene Zeiten, da Sinne und Intellect nur Material sammelten, ohne Prüfung, Sichtung, Anordnung, ohne Gleiches zu Gleichem zu legen und die widersprechenden Einzelfälle zu eliminiren. Das abstracte Denken ist noch unentwickelt, das Gehirn kann nicht arbeiten, ohne dass anschauliche Bilder auftauchten und sich mit bewegten. Alles in allem ist der Traumzustand, wie jener Culturzustand, höchst schöpferisch und intensiv, aber ohne Gesetz; der Instinct für Werth oder Unwerth, für Bevorzugung, Annahme oder Ausschluss der einzelnen Bilder ist ausgeschaltet: es giebt keine Jury und „Hängecommission“, sondern Jeder stellt aus, was er gemalt hat, und hängt es hin, wo er Platz findet.
\par

